\chapter{Existence and uniqueness of solutions}

\prerequisites{linear algebra (isomorphism, injectivity, surjectivity of linear maps), a brief knowledge on Banach spaces, continuity, properties of continuous functions on compact sets, brief knowledge on ordinary differential equations.}

If you're reading this far, I suppose you are already comfortable with calculus, the
methods of solving differential equations, some linear algebra, and perhaps some
real analysis. This chapter concerns the general theory of differential equations,
consisting mostly of the analysis of existence and uniqueness of solutions. It is not
about how to solve them, or how they're applied in mechanics. For those of the applied
side, refer to the first book on foundations of calculus. But I greatly welcome those
who are interested as this chapter is going to be a very fulfilling journey that lies
in the intersection between physics, linear algebra, analysis, and some very cute visualizations.

The power of classical mechanics comes from the ability to \emph{accurately} predict the
future and past of a physical system given enough information on the current configuration
of a system. E.g., consider an oscillating spring system in a world with drag and no external
disturbances. Naturally, it should gradually lose power of oscillation over time; the
drag would make the mass lose energy until it stops at \(t \appr \infty\). And in the
past, the spring must've oscillated more wildly as the system hasn't lose energy to the drag
yet. Prior works in calculus also have shown that the Newton's second law governing this system,
\begin{equation}
  m \odv[ord = 2]{x}{t} = - b \odv{x}{t} - kx
\end{equation}
has the solution
\begin{equation}
  x(t) = A \e^{-\beta t} \cos\ab(\sqrt{\omega_0^2 - \beta^2} t - \phi)
\end{equation}
where the cosine represents the oscillation, and the exponential term represents the decay.

It'd be very na\"ive to assume that this determinism works with every system. E.g., consider
a leaking water bucket that's shaped like a cone. Eventually, it loses all its water.
But if I give you the state of the system after the bucket has lost all its water, can
you trace back when the last drop of water leaked out? Of course not! No matter where
in time the last drop has oozed out, you're going to be left with a persistently empty
bucket, assuming that no one trolls you and put more water in a bucket. So, determinacy
is lost here: there is no laws which can reverse this process.

Another system that has this property would be the falling object on the ground. After
the object has landed and rested still, can you trace back when exactly the object touches
the ground? It could've been one second ago or one year ago! If two objects are dropped from the
exact same height, it would land and rest in the same manner, reaching a state of persistence rest.
When this state is reached, the information about when the object was dropped is completely lost.

Two of these systems I gave have one thing in common: the lost of something causing time
to not be reversible --- water losing from the bucket into oblivion, the ball losing
the energy into the ground. And you might think: if you managed to trace all those lost
--- if you know where the water goes and how it pools the ground, or how the ground
dissipates those energy from the balls into heat, etc., you would \emph{eventually} be
able to determine the object's complete history. And that is true in a sense. However,
John Norton has a counterexample to this ideology: Norton's dome \cite{norton-dome},
in which I'd like to state a modified version of Norton's dome in two dimensions instead.

Norton's dome is described by the equation
\begin{equation}
  y = - \frac{2}{3g} r^{\frac{3}{2}}
\end{equation}
where \(r\) is the arclength along the dome. Initially, we place a ball of mass \(m\)
resting on the top of the dome with no velocity nor acceleration. The acceleration of
the part that's tangent to the dome is then
\begin{equation}
  \odv[ord = 2]{r}{t} = b^2 \sqrt{r}.
\end{equation}
in which there are actually two family of solutions to this differential equation:
\begin{equation}
  r(t) = 0, \mathand r(t)
       = \begin{cases}
           0                      & \textrm{if}~ t \leq T \\
           \frac{1}{144}(t - T)^4 & \textrm{if} ~ t > T
         \end{cases}
\end{equation}
where \(T \in \mathbb{R}\). Hopefully it is easy enough to check that these two solutions
actually satisfies the Newton's equation by simply taking the derivative.

The first family is what we expect: the ball just stays forever rested on the dome
without falling off. But in the second family, the ball somehow manages to roll off
the dome at any time \(T\) without any outside forces: \(r(t)\), the distance that the
ball rolled, is positive after \(t > T\) where \(T\) is \emph{arbitrary}.

The question then becomes: what solution do we trust, and is the Newton's law wrong,
or is Newton's law only applicable on some edge cases. This is a problem of philosophy
of physics which we shall not touch. For the interested readers, one can refer to
\cite{malament-2008}, or a simple Google Search would bring up all the heated debates
around this particular physical systems. The philosophy in question is greatly beyond
the scope of this book, so we shall stroll back a step and answer a more mathematically
approachable question: when is Newton's second law not deterministic. Or more generally,
\begin{displayquote}
  \emph{When does a differential equation possess a single solution?}
\end{displayquote}

This question, or rather this scheme of thinking, is one part of core mathematics
called \textbf{well-posedness}. Mathematicians, at least I, tend to check first if the
problem is \textbf{well-posed}, or solvable in the first place before actually dip our
hands in and solve it. A problem is well-posed if
\begin{enumerate}
  \item There exists a solution, and
  \item The solution is unique. I.e., if there is a solution, there must be only one solution
  \item The behavior of the solution changes smoothly with the auxilary conditions.
\end{enumerate}
In the next section, we shall explore these behaviors via visualizing slope fields.

\section{A quick recap on definitions of differential equations and phase space}

\begin{df}{ODE of one variable}{}
  A differential equation is an equation of the form
  \begin{equation}
    f(x, y, y', y'', \dots, y^{'(n)}) = 0
  \end{equation}
  where the derivatives are understood to be w.r.t. \(x\). The order of the differential
  equation is then the highest order derivative that appears.

  A function \(\phi(x)\) is a solution of the differential equation if when \(y\) is
  replaced with \(\phi\), \(f(x, \phi, \phi', \dots, \phi^{'(n)}) = 0\).
\end{df}

This definition, even though a very precise one that captures all the differential
equations you could ever think of, is albeit, useless as the family is too large for
us to analyze any meaningful behaviors. Hence in this chapter, we shall start with the
easiest one: first order differential equations with a specific form here:

\begin{df}{Standard first order ODE of one variable}
  A first order differential equation in standard form can be written as
  \begin{equation}
    \odv{y}{x} = f(x, y).
  \end{equation}
\end{df}

Equations of this form are particularly nice, as there is a clear interpretation of it
which I'd like you to follow along.

Consider any arbitrary functions \(f(x, y)\)

\section{Existence and uniquness theorem - Euler's method}

The theorem that we shall derive in this section is the Peano's existence theorem and
Cauchy's existence theorem which concerns ordinary differential equations of the form
\(\odv{y}{x} = f(x, y(x))\). Here, I heavily focus on the word \emph{derive}, as I'll
guide you through and motivate the entire thought process of solving this problem. In
most courses, the theorem will just be handed straight to you and you would be confused
on why such requirements affect the solutions of the differential equations at all. Other
than that, this problem is just tricky. How would you even go about proving whether a
differential equation has a solution or not when you don't even know what function you're
dealing with. Even if you do, sometimes a function might even be impossible to explicitly integrate.
So, hopefully at the end of this section, every bits and pieces will begin to make sense again.

Since we would like to prove that a solution \emph{exists}, an approach would be to construct
the solution itself: to approach the problem constructively. Of course, there is no way
to explicitly construct such solution, but there is a way to approximate it that's thoroughly
developed in the field of scientific computing and numerical analysis. And gladly enough,
they're also interested in whether the solution is correct and is stable or not. So,
we can use their analysis of \emph{correctness} to motivate our uniqueness and existence theorem.

The method of choice here will be the Euler's method which is used to find the \emph{next
step} of the solution. Say we have an initial value problem
\begin{equation}
  \odv{y}{x} = f(x, y(x)), \mathand y(x_0) = y_0
\end{equation}
on some interval \(I \in \mathbb{R}\). Preferably, \(I\) should be a connected interval
that extends from \(x_0\) to \(x_0 + X\) for some positive \(X\) so that we can try to
find what points follows \(y(x_0)\) along the solution curve. For that, let's try to
find \(y(x_0 + \epsilon)\) for some very small \(\epsilon\). Note that if we want to
extend \(y(x_0)\) to the left, we can perform a mirror coordinate transformation around \(x_0\).

Here's a quick recap of the Euler's method. First, rewrite the differential equation
as a discrete equation:
\begin{align}
  \odv{y}{x}_{x_0} \approx \frac{y(x_0 + \epsilon) - y(x_0)}{\epsilon} & = f(x_0, y(x_0))                    \\
  y(x_0 + \epsilon)                                                    & = \epsilon f(x_0, y(x_0)) + y(x_0).
\end{align}
Also as one can simply induct,
\begin{equation}
  y(x_0 + \alpha\epsilon) = \epsilon\ab[\sum_{k = 0}^{\alpha - 1} f\ab\big(x_0 + k\epsilon, y(x_0 + k\epsilon))]
                            + y(x_0). \label{eq:existenceuniqueness-euler}
\end{equation}

We can then choose our step size \(\epsilon\) to partition the interval \(x_0\) to
\(x_0 + X\) into equal sections, i.e., let \(N\epsilon = X\) for some \(N \in \mathbb{N}\).
We can then construct a sequence
\begin{equation}
  y(x_0), y(x_0 + \epsilon), y(x_0 + 2\epsilon), y(x_0 + 3\epsilon), \dots, y(x_0 + N\epsilon).
\end{equation}
In order to form this sequence into a proper approximation, we can join the dots using
straight lines of constant slope. Our approximated solution \(\phi_N(x)\) is then a
piecewise function. For \(x \in \cintv{x_0 + n\epsilon, x_0 + (n + 1)\epsilon}\),
\begin{align}
  \phi_N(x) & = \begin{multlined}[t]
                  y(x_0 + n\epsilon) \\
                  + \frac{y(x_0 + (n + 1)\epsilon) - y(x_0 + n\epsilon)}{\epsilon}\ab\big(x
                  - (x_0 + n\epsilon))
                \end{multlined} \\
            & = y(x_0 + n\epsilon) + f\ab\big(x_0 + n\epsilon, y(x_0 + n\epsilon))\ab\big(x - (x_0 + n\epsilon)). \label{eq:existenceuniqueness-euler_phi}
\end{align}
Or, \(\phi_N(x)\) can be specified by piecewise function of constant derivatives:
\begin{equation}
  \odv{\phi_N}{x}_{x_0} = f\ab\big(x_0 + n\epsilon, y(x_0 + n\epsilon))(x - (x_0 + n\epsilon)), \quad \phi_N(x_0 + n\epsilon)
                        = y(x_0 + n\epsilon).
\end{equation}
We then construct the sequence of functions \(\sequence{\phi_N}{N = 1}{\infty}\). If this
sequence behaves well enough in the sense that it uniformly converges to some function, then that
function is our solution, and we would've shown that a solution exists. But under what condition?

Firstly, the solution to the differential equation shouldn't have arbitrarily large
vertical oscillations. Otherwise, Euler's method would blow up. From
\begin{equation}
  y(x_0 + \epsilon) - y(x_0) = \epsilon f(x_0, y(x_0)), \nonumber
\end{equation}
this sum has to be strictly finite for the solution to be physical. Since \(\epsilon\)
can be any number, \(f(x_0, y(x_0))\) must remain finite on the domain, i.e., bounded.

However, notice that \(f\) has two arguments: \(x_0\) and \(y(x_0)\). Specifying
boundedness on only the first argument wouldn't be enough to guarantee that nothing
would blow up. Hence, we must require \(f\) to be bounded on both \(x\) and \(y\)
axis. Specifically, on \(I_x \times I_y\) where \(I_x, I_y \subseteq \mathbb{R}\) are
intervals. Intuitively, this draws a rectangle in \(\mathbb{R}^2\)

Now let's analyze what happens when \(\epsilon\) shrinks. For simplicity, let
\(\epsilon' = \frac{\epsilon}{2}\). Then, we should intuitively get
\begin{equation}
  y(x_0 + 2\epsilon') \approx y(x_0 + \epsilon)
\end{equation}
where the L.H.S. is evaluated by using the Euler's method twice with a step of
\(\epsilon'\), and the R.H.S. steps only once. This approximation can be restated as
\begin{equation}
  |y(x_0 + 2\epsilon') - y(x_0 + \epsilon)| < \delta
\end{equation}
where you can always select \(\epsilon\) in a way such that \(\delta\) can get arbitrarily
small. Expanding the approximation above by \cref{eq:existenceuniqueness-euler}, we get
\begin{align}
  y(x_0) + \epsilon'\ab\Big[f\ab\big(x_0 + \epsilon, y(x_0 + \epsilon)) + f\ab\big(x_0, y(x_0))] & \approx y(x_0) + \epsilon f(x_0, y(x_0)) \nonumber \\
  \frac{\epsilon}{2}\ab\Big[f\ab\big(x_0 + \epsilon, y(x_0 + \epsilon)) + f\ab\big(x_0, y(x_0))] & \approx \epsilon f(x_0, y(x_0))                    \\
  f\ab\big(x_0 + \epsilon, y(x_0 + \epsilon)) + f\ab\big(x_0, y(x_0))                            & \approx 2 f(x_0, y(x_0))                           \\
  f\ab\big(x_0 + \epsilon, y(x_0 + \epsilon))                                                    & \approx f\ab\big(x_0, y(x_0))
\end{align}
or that there exists some \(\epsilon\) in which
\begin{equation}
  |f\ab\big(x_0 + \epsilon, y(x_0 + \epsilon)) - f\ab\big(x_0, y(x_0))| < \delta
\end{equation}
for all \(\delta\). This is the exact requirement for continuity: it is required for
the Euler's equation to give \emph{consistent result} for all \(\epsilon\). In the
following proof, we shall see how these two conditions --- boundedness and continuity
--- give rise to a solution for the differential equation.

These two conditions are enforced by requiring that \(I = I_x \times I_y\) is a closed
rectangle in \(\mathbb{R}^2\). By the Heine-Borel theorem, a closed set on \(\mathbb{R}^n\)
is compact, and any continuous function on a compact set is automatically bounded.

Before we move on to the formal proof, let's see some examples where an initial value
problem just refuses to have any solution. Note that the examples may look a bit exotic.
But that's just because we're only requiring \(f\) to be continuous and bounded on
\(I\) which already covers most of the function that is normally seen.

{\color{red} W.I.P. illustrate some examples on Euler's equation}

Now, to the formal proof of the theorem.

\begin{theorem}{Peano's existence theorem}{}
  For a differential equation \(\odv{y}{x} = f(x, y(x))\) where \(f\) is continuous
  on a compact set \(\mathcal{I} \subset \mathbb{R}^2\), an initial value problem with
  \(y(x_0) = y_0\) where \((x_0, y_0) \in I\) has a solution.
\end{theorem}

\begin{proof}
  Consider a subset of \(\mathcal{I}\),
  \begin{equation}
    I = \cintv{x_0 - a, x_0 + a} \times \cintv{y_0 - b, y_0 + b}
  \end{equation}
  This subset represents a compact rectangle in two dimensions. Since \(f\) is required
  to be continuous on a compact set, it is bounded by some real number \(M\).

  Let's now construct the solution on the domain interval \(I_x = \cintv{x_0, x_0 + X}\)
  where \(X = \min\ab(a, \frac{b}{M})\). The reason for this choice of interval is that
  in some cases, the function might be too steep that it actually shoots passes the
  upper wall of the rectangle \(I\). Hence, we restrict the domain of \(I\) to assure
  that for all \(x \in \cintv{x_0, x_0 + X}\), \((x, \phi_n(x)) \in I\).

  Partition \(I_x\) into equal intervals of \(x_{0, k}, x_{1, k}, \dots, x_{n, k}\) with
  \begin{equation}
    x_{n, k+1} - x_{n, k} = \frac{X}{n} = h_n.
  \end{equation}

  Construct the sequence of approximations of functions \(\sequence{\phi_n(x)}{n = 1}{\infty}\)
  as shown earlier in \cref{eq:existenceuniqueness-euler_phi}. First, we prove that
  this sequence doesn't leave \(I\).

  For all \(x \in I_x\),
  \begin{equation}
    |\phi_n(x) - y_0| \leq M(x - x_0) \leq MX \leq b;
  \end{equation}
  hence, the graph is inside \(I\) for all \(x \in I_x\) and \(n \in \mathbb{N}\).

  Then, we show that this sequence is uniformly bounded and is equicontinuous in order
  to apply Arzel\'a-Ascoli theorem, yielding a uniformly convergent subsequence.

  Since \(|\phi_n(x)| \leq |y_0| + b \in I_y\) for all \(x \in I_x\) and \(n \in \mathbb{N}\),
  \(\phi_n(x)\) is uniformly bounded.

  Now for all \(x_1, x_2 \in I_x\),
  \begin{equation}
    |\phi_n(x_1) - \phi_n(x_2)| \leq M|x_1 - x_2|;
  \end{equation}
  hence, \(\phi_n\) is Lipschitz with Lipschitz constant \(M\) throughout all \(n\).
  This family is then equicontinuous.

  By the Arzel\'a-Ascoli theorem, there is then a subsequence of functions,
  \(\sequence{\phi_{n_k}(x)}{n = 1}{\infty}\), that uniformly converges.

  For ease of notation, assign for each \(n\), a function \(P_n(x) = x_{n, k}\) which
  denotes the starting point of the segments of the partition that \(x\) is in. Then,
  for all points that doesn't belong in the partition,
  \begin{equation}
    \odv{\phi_n}{x} = f\ab\big(P_n(x), \phi_n(P_n(x)))
  \end{equation}
  Now, define the error of the approximation's derivative to be
  \begin{align}
    E_n(x) & = \odv{\phi_n}{x} - f(x, \phi_n(x))                          \\
           & = f\ab\big(P_n(x), \phi_n(P_n(x))) - f\ab\big(x, \phi_n(x)).
  \end{align}

  Since \(I\) is compact, and \(f\) is continuous, it is also \(f\) uniformly continuous. I.e.,
  \begin{multline}
    \forall (\epsilon \in \mathbb{R}_{>0}) \exists (\delta \in \mathbb{R}_{>0}) \forall ((x_1, y_1), (x_2, y_2) \in I) : \\
    \ab(|x_1 - x_2| < \delta \land |y_1 - y_2| < \delta) \implies |f(x_1, y_1) - f(x_2, y_2)| < \epsilon.
  \end{multline}

  By the presumed continuity, this error termed can be bounded by the distance between
  \((x, \phi_n(x))\) and \((P_n(x), \phi_n(P_n(x)))\). Along the first axis,
  \begin{equation}
    |x - P_n(x)| < h_n.
  \end{equation}
  Along the second axis, we can use the mean value theorem to say that there exists
  some \(c\) in the segment such that
  \begin{equation}
    |\phi_n(x) - \phi_n(P_n(x))| = \ab|\odv{\phi_n}{x}_c||x - P_n(x)| \leq M h_n
  \end{equation}

  If we let \(h_n = \min\ab(\epsilon, \frac{\epsilon}{M})\), rounded down to the lowest
  that still maintains equipartition, it follows that
  \begin{equation}
    |x - P_n(x)| < \epsilon \land |\phi_n(x) - \phi_n(P_n(x))|.
  \end{equation}
  From uniform continuity of \(f\), we also get that
  \begin{equation}
    |E_n(x)| = |f\ab\big(P_n(x), \phi_n(P_n(x))) - f(x, \phi_n(x))| \leq \epsilon;
  \end{equation}
  hence, \(|E_n(x)| \appr 0\) uniformly aswell.

  The last step is to verify that the limit of this function subsequence actually
  satisfies the differential equation.

  Let \(x \in I_x\), and choose a step size \(h \in \mathbb{R}_{>0}\) s.t. \(h < X\).

  By the continuity of and piecewise linearity of \(\phi_n(x)\), the difference
  \(\phi_n(x + h) - \phi_n(x)\) is just \(h\) times the weighted average of the
  slopes on each segments. By the intermediate value theorem, there exists a point
  \(\tau_n \in \cintv{x, x + h}\) such that
  \begin{align}
    \frac{\phi_n(x + h) - \phi_n(x)}{h} & = f\ab\big(P_n(\tau_n) \phi(P_n(\tau_n)))         \\
                                        & = f\ab\big(\tau_n, \phi_n(\tau_n)) + E_n(\tau_n).
  \end{align}
  Now take the limit of the expression above by letting \(n = n_k\) and take \(k \appr \infty\).

  On the L.H.S.,
  \begin{equation}
    \lim_{k \appr \infty} \frac{\phi_{n_k}(x + h) - \phi_{n_k}(x)}{h} = \frac{y(x + h) - y(x)}{h}
  \end{equation}
  because \(\phi_{n_k}\) connverges uniformly to the solution \(y\). On the R.H.S.,
  \begin{equation}
    \lim_{k \appr \infty} f(\tau_{n_k}, \phi_{n_k}(\tau_{n_k})) + E_{n_k}(\tau_{n_k}) = f(\tau, y(\tau))
  \end{equation}
  for some \(\tau\) because \(E_{n_k}\) uniformly converges to zero. Hence,
  \begin{equation}
    \frac{y(x + h) - y(x)}{h} = f\ab\big(\tau, y(\tau)).
  \end{equation}
  Then, take the limit as \(h \appr 0^+\). On the R.H.S., since \(\tau \in \cintv{x, x + h}\),
  \(\tau \appr x\) as \(h \appr 0^+\). We then get
  \begin{equation}
    \odv{y}{x} = \lim_{h \appr 0^+}\frac{y(x + h) - y(x)}{h}
               = \lim_{\tau \appr x} f(\tau, y(\tau))
               = f(x, y(x)).
  \end{equation}
  Since \(y(x_0) = \lim_{k \appr \infty} \phi_{n_k}(x_0) = y_0\), \(y(x)\) satisfies
  \(y'(x) = f\ab\big(x, y(x))\) and \(y(x_0) = y_0\) on \(I_x\).

  This proves that the sequence of approximations earlier always uniformly converges
  to a valid solution, showing that the solution must exists. For the direction
  \(\cintv{x_0 - X, x_0}\), perform a mirror substitution.
\end{proof}

This already assures us that most differential equations that we're interested in has
a solution, but that's only a part of a well-posed problem. The other is uniqueness. This
is important in deterministic theories where we want there to be only \emph{one single trajectory}
that a system can take. So, what conditions do we have to further impose to yield determinism?

When we constructed a sequence of Euler approximations, we used Arzel\'a-Ascoli theorem to
show that there exists a uniformly convergent subsequence. However, it does not guarantee that there
is only \emph{one} subsequence. Different subsequences might converge to a different solutions.
So to fix that, we require the \emph{entire sequence} of \(\phi_n\) to uniformly converge.

So, let's consider any arbitrary Euler's approximation:
\begin{equation}
  \phi_n(x_{n, k+1}) = \phi_n(x_{n, k}) + h_n f\ab\big(x_{n, k}, \phi_n(x_{n, k})).
\end{equation}
We might try to bound this \emph{directly} to the real solution \(y(x)\). Since \(y\)
is differentiable, the mean value theorem applies to give
\begin{equation}
  y(x_{n, k+1}) = y(x_{n, k}) + h_n f\ab\big(\eta_k, y(\eta_k))
\end{equation}
for some \(\eta \in \cintv{x_{n, k}, x_{n, k + 1}}\). Subtracting yields
% \begin{equation}
%   \phi_n(x_{n, k + 1}) - y(x_{n, k + 1}) 
% \end{equation}

\section{Existence and uniqueness theorem - integral equations}

The proof earlier using the Euler's method, even with its clear intuition on the
constraints of functions, can be clunky because we're trying to interpolate discrete
points into approximating functions. But what if, we start with a smooth function to begin with?

Notice that the differential equation \(\odv{y}{x} = f(x, y(x))\) can also be casted
into an integral form:
\begin{align}
  \int_{x_0}^{x} \odv{y}{x'} \odif{x'} & = \int_{x_0}^x f(x', y(x')) \odif{x'}           \\
  y(x) - y(x_0)                        & = \int_{x_0}^x f(x', y(x')) \odif{x'}           \\
  y(x)                                 & = \int_{x_0}^x f(x', y(x')) \odif{x'} - y(x_0).
\end{align}
This integral form is actually equivalent to the original differential equation, and
we get the boundedness of \(f(x', y(x'))\) just from looking at the integral --- an
integrable function must be bounded! However, the continuity is a bit unclear

\section{Linear differential equations}

Here we treat a special class of differential equations: linear equations. As we might've
seen in earlier chapters on linear algebra, the derivative \emph{is} a linear operator in
the sense that if \(f\) and \(g\) are differentiable functions and \(a\) is a scalar constant, then
\begin{equation}
  \odv*{\ab(a f(x) + g(x))}{x} = a \odv{f(x)}{x} + \odv{g(x)}{x}
\end{equation}
Hence, it would be perfectly reasonable to replace the Liebniz' derivative notation
with the operator \(D\). A differential equation then takes on the form
\begin{equation}
  F(t, y, Dy, D^2 y, \dots, D^n y) = 0.
\end{equation}
A linear differential equation is then defined as follows:

\begin{df}{Linear differential equation}{}
  A linear differential equation on a domain \(I \subseteq \mathbb{R}\) is an equation of the form
  \begin{equation}
    a_n(t) D^n y + a_{n-1}(t) D^{n-1} y + \dots + a_1(t) Dy + a_0(t)y = g(t)
  \end{equation}
  where \(a_i(t)\) are all defined on \(I\).

  A function \(\phi(t)\) is a solution of the differential equation above \emph{iff}
  for all \(t \in I\),
  \begin{equation}
    a_n(t) D^n \phi + \dots + a_0(t) \phi = g(t)
  \end{equation}
\end{df}
Since this operator is linear, we can factor out the differential operator and write
\begin{equation}
  \ab(a_n(t) D^n + \dots + a_1(t) D + \iden) y = g(t)
\end{equation}
or just
\begin{equation}
  P(D, t)y = g(t)
\end{equation}
where \(P(z, t)\) is a polynomial of the differential operator defined by
\begin{equation}
  P(z, t) = a_n(t) z^n + \dots + a_1(t) z^1 + a_0(t):
\end{equation}
The problem of solving the differential equation then boils down to finding the inverse
of \(P(D, t)\). However, we run to a problem:

\begin{prop}{}{}
  Let \(I \in \mathbb{R}\) be an interval. The differential operator
  \(D: C^1(I, \mathbb{R}) \to C^0(I, \mathbb{R})\) is surjective but not injective.
\end{prop}

This lemma is analogous to why we must add a constant term (\(+ C\)) when we're taking
the antiderivative, which is the "inverse" of the differential operator.

\begin{proof}
  Pick a point \(c \in I\), then for all \(g \in C^0(I, \mathbb{R})\),
  \begin{equation}
    f(x) = \int_c^x g(t) \odif{t}
  \end{equation}
  is in \(C^1(I, \mathbb{R})\) and \(Df = g\) by the fundamental theorem of calculus.
  Hence, \(D\) is surjective.

  Let \(f \in C^1(I, \mathbb{R})\) and \(c\) be the constant function. If \(Df = g\),
  then \(D(f + c) = Df + Dc = Df\): \(D\) is not injective.
\end{proof}

Hence, we can only ask for a mere preimage of \(D\). Here, we can harbour some hope from
the property of linear operators themselves. First, the solution of a non-injective linear equation.

\begin{lemma}{Solution of a non-injective linear equation}{}
  If \(L \in \mathcal{L}(V, W)\) and \(y \in W\), then every solution of \(Ly = x\) is
  of the form \(y_0 + \ker L\).
\end{lemma}

Note that this proposition doesn't mention anything about differential equations. It is just true
to the nature of \emph{any} linear equations. Differential equations are a special case of this.

\begin{proof}
  First, to show that \(\phi = y_0 + \xi\) where \(\xi \in \ker L\) is a solution:
  \begin{equation}
    L\phi = L(y_0 + \xi) = Ly_0 + L\xi = Ly_0 = x.
  \end{equation}

  Then, to show that every solution is of this form, let \(\psi\) be any other solution
  to \(Ly = x\). Then,
  \begin{equation}
    L(\phi - y_0) = L\phi - Ly_0 = 0,
  \end{equation}
  i.e., \(\phi - y_0 \in \ker L\) and is equal to some \(\xi \in \ker L\). Rearranging, we get
  \begin{equation}
    \psi = y_0 + \xi:
  \end{equation}
  concluding the proof.
\end{proof}

This trivially yields

\begin{prop}{}{}
  For a linear differential equation \(P(D, t)y = g(t)\), if \(\phi(t)\) is a solution
  and \(P(D, t) \psi(t) = 0\), i.e., \(\psi(t) \in \ker P(D, t)\), then \(\psi(t) + \phi(t)\)
  is also a solution.
\end{prop}

\begin{proof}
  Follows directly from the linearity of \(P(D, t)\) and above lemma.
\end{proof}

This is exactly why when we solve linear differential equations, we must consider both
the original equation itself and its auxilary equation \(P(D)y = 0\) as the complete set of solution
is always the sum of the \textbf{particular solution} of the original equation and the kernel of
\(P(D)\). The solution in the kernel of \(P(D)\) is also called the \textbf{complementary solution}.

Note that the particular solution here need not be unique. If \(\psi(x)\) and \(\phi(x)\)
are solutions, then they differ by \(\xi(x)\) where \(\xi \in \ker P(D)\). One can write
\begin{equation}
  \phi(x) = \psi(x) - \xi(x), \mathor \psi(x) = \phi(x) + \xi(x).
\end{equation}

A bit of pondering should arise a series of questions. When we solve differential
equations (not the initial value problem), we often put \(+ C\) in order to keep track
of the \emph{complete set of solutions}. I.e., we're saying that the set
\begin{equation}
  \{f(x, C_1, \dots, C_n) : C_1, \dots, C_n \in \mathbb{R}\}
\end{equation}
forms the entire solution set, and no more solution lies outside this set. The amount
of \(C_i\)'s in this function represents some kind of degree of freedom in our solution
space. And by intuition, an \(n\)-th order differential equation \emph{should} have \(n\)
degrees of freedom, as every new constant \(C_i\) added represents an integration; an
\(n\)-th order differential equation has to be integrated \(n\) times.

Generally, nothing more can be said about the nature of the solution set. But with
linear equations, we can. As the solution set is governed by \(y_0 + \ker P(D)\), all
of the degrees of freedom must reside in \(\ker P(D)\). Since the kernel of any linear
operator is a vector space, it has a well-defined dimension (given that it's not
infinite-dimensional). This dimension number directly corresponds to the degree of freedom
in the solution set. Meaning that, if \(P(z)\) is of degree \(n\), \(\ker P(D)\) must
also have dimension \(n\). The following proposition states this in more detail:

\begin{prop}{Dimension of the solution set of linear equations}{}
  For a linear differential equation \(P(D, t)y(t) = 0\) on a closed and bounded interval \(I\)
  where the coefficients of the highest order derivative is non-zero, \(\dim \ker P(D) = \deg P\).
\end{prop}

\begin{proof}
  We shall show this through methods of initial value problems. I.e., we show that there exists
  an isomorphism between the initial condition of the differential equations and \(\ker P(D)\).

  As a bit of simplification, since the coefficients of the highest order derivative
  is non-zero in \(I\), we can divide through to yield another differential equation
  \begin{equation}
    Q(D)y = 0 \quad \textrm{where} \quad Q(D) = \frac{P(D)}{a_n(t)}
  \end{equation}
  that has the same solution set as \(P(D)y = 0\).

  Note that by the equation, \(y(t) \in \ker P(D)\).

  Fix some \(t_0 \in I\) and define a map \(T : \ker P(D, t) \to \mathbb{R}^n\) by
  \begin{equation}
    Ty = \ab(y, \odv{y}{t}_{t_0}, \odv[ord = 2]{y}{t}_{t_0}, \dots, \odv[ord = n - 1]{y}{t}_{t_0}).
  \end{equation}

  \(T\) is definitely linear, as the derivatives of a function are linear.

  Now to prove injectivity, suppose \(Ty = 0\), then
  \begin{equation}
    y(t_0) = \odv{y}{t}_{t_0} = \dots = \odv[ord = n - 1]{y}{t}_{t_0} = 0.
  \end{equation}
  This initial value problem has a unique solution because
  \(\frac{a_i}{a_n} \in C^1(I, \mathbb{R})\) on a compact interval, i.e., all
  coefficients of \(Q(D)\) are Lipschitz. Since \(0\) is one of the solution as \(0\)
  is in the kernel of every linear maps, it is the only solution and \(T\) is injective.

  To prove surjectivity, assume that \(Ty = Ty' = (s_1, \dots, s_n)\). Since the initial
  values are equal, they must correspond to the same unique solution by the uniqueness
  theorem. Hence, \(y = y'\) and \(T\) is surjective.

  Since \(T\) is linear, injective, and surjective, it is an isomorphism. Two isomorphic
  vector space always have the same dimension. Hence, \(\dim \ker P(D) = n = \deg P(z)\).
\end{proof}

One very big consequence of the lemma above is: if you managed to find \(n\) solutions
that are linearly independent, then the span of those solutions is the entire set. This
is a basic consequence of linear algebra as a linearly independent list of the right length always
form a basis. Hence, we shall formally introduce the notion of linear independence of functions.

\begin{df}{Linear independence of functions}{}
  A list of functions \(\{f_1(t), \dots, f_n(t)\}\) are linearly independent \ifft for
  all \(t\) in their natural domains, the only way to make
  \begin{equation}
    c_1 f_1(t) + \dots + c_n f_n(t) = 0
  \end{equation}
  is for \(c_1 = \dots = c_n = 0\). The said list is linearly dependent otherwise.
\end{df}

One way to compute whether a set of functions are linearly independent or not is to
explicitly provide a proof. E.g.,

\begin{prop}{The set of polynomials of different degrees are linearly independent}
  The set \(\{x^{k_1}, x^{k_2}, x^{k_3}, \dots, x^{k_n}\}\) where \(k_1 \neq \dots \neq k_n\)
  is linearly independent.
\end{prop}

\begin{proof}
  Let's show that \(\{x^1, \dots, x^n\}\) are linearly independent. A subset of a linearly
  independent set is also linearly independent hence any set \(\{x^{k_1}, \dots, x^{k_n}\}\)
  must also be a linearly independent set.
\end{proof}

{\color{red} W.I.P. Linear independence of polynomials, Wronskians, Abel's theorem}

To formalize the method of undetermined coefficients discussed earlier, there is quite
a useful corollary:
\begin{cor}{}{}
  Polynomials of the derivative operator with constant coefficients are commutative
\end{cor}
\begin{proof}
  Let \(P(D) = \sum_{i=0}^n p_i D^i\) and \(Q(D) = \sum_{i=0}^m q_i D^i\) be polynomials
  of the derivative operator with constant coefficients
  \begin{align}
    P(D)Q(D) & = \sum_{i=0}^n p_i D^i \sum_{j=0}^m q_j D^j                     \\
             & = \sum_{i=0}^n \sum_{j=0}^m p_i D^i q_j D^j                     \\
             & = \sum_{i=0}^n \sum_{j=0}^m q_j D^j p_i D^i                     \\
             & = \sum_{i=0}^m q_i D^i \sum_{j=0}^n p_j D^j = Q(D)P(D) \qedhere
  \end{align}
\end{proof}

Note that if the coefficients are non-constant, aka variable, the commutativity property is lost.

\begin{theorem}{Solution of homogeneous linear differential equations with constant coefficients}{}
  For the linear differential equation \(P(D)y = 0\) with
  \begin{equation}
    P(D) = \prod_{i=0}^{n} (D - \lambda_i \iden)^{k_i},
  \end{equation}
  the solution set is then completely spanned by
  \begin{equation}
    \{ x^r \e^{\lambda_i x} : (r \in \mathbb{N}_{<k}) \land (i \in \mathbb{N}_{\leq n}) \}
  \end{equation}
\end{theorem}

\begin{proof}
  We first prove that the given basis is a part of the solution set, then prove that
  every solution can be written as a linear combination of these basis.

  Consider a more general function \(f(x)\e^{\lambda x}\).
  \begin{align}
    (D - \lambda \iden) [f(x) \e^{\lambda x}] & = f D \e^{\lambda x} + \e^{\lambda x} Df - \lambda f(x) \e^{\lambda x} \\
                                              & = \e^{\lambda x} Df
  \end{align}
  By simple induction which is left as an exercise,
  \begin{equation}
    (D - \lambda \iden)^k [f(x) \e^{\lambda x}] = \e^{\lambda x} D^k f
  \end{equation}
  If we want the function of this form to be a solution, i.e.,
  \(f(x)\e^{\lambda x} \in \ker (D - \lambda \iden)^k\), then we can set \(f(x) = x^r\)
  where \(r \in \mathbb{N}_{<k}\) and form a set
  \begin{equation}
    S_i = \{x^r \e^{\lambda_i x} : r \in \mathbb{N}_{<k_i}\}.
  \end{equation}
  Note that this set is linearly independent. The span of this set is then also a
  subset of \(\ker (D - \lambda_i \iden)^{k_i}\) which has dimension \(k_i\). Since
  \(\deg (D - \lambda_i \iden)^{k_i} = k_i\), this spanning set is of the correct length
  hence constitute a basis for \(\ker (D - \lambda_i \iden)^{k_i}\).

  By the commutativity of polynomials of constant coefficients,
  \begin{equation}
    \ab[\prod_{i = 0}^n (D - \lambda_i \iden)^{k_i}] y = \ab[\prod_{i = 0, i \neq j}^n (D - \lambda_i \iden)^{k_i}] (D - \lambda_j \iden)^{k_j} y
                                                       = 0.
  \end{equation}
  Hence, \(y\) can be in \(\ker (D - \lambda_j \iden)^{k_j}\) for all \(j \in \mathbb{N}_{<n}\).
  Each one of these form its own set \(S_j\).

  It's obvious that
  \begin{equation}
    \mathcal{S} = \bigcup_{j=0}^{n} S_j
  \end{equation}
  is a linearly independent set with cardinality \(\sum_{i=0}^n k_i\). Hence, we can
  construct the span of \(\mathcal{S}\) with each elements \(s(x)\) in the span satisfying
  \(P(D)s(x) = 0\). The span of \(\mathcal{S}\) is then the entirety of \(\ker P(D)\) because
  \begin{equation}
    \deg P(D) = \dim \spn \mathcal{S} = \sum_{i=0}^n k_i.
  \end{equation}
\end{proof}

Note that the theorem above does not care about whether each of the \(\lambda_i\) are
complex or reals. For complex factorizations, one can use the Euler's formula to expand
\(\e^{\lambda_i x}\) into cosines and sines. Suppose \(\lambda_i = a + b\im\), then
\begin{equation}
  \e^{a+b\im} = \e^a \e^{b\im} = \e^a (\cos b + \im \sin b).
\end{equation}
However, note that there are nothing that assures that the solution is real. The realness
of the solution only comes from the configuration of the system on the initial value
problem. But for second order linear differential equations, the result is guaranteed
to be real by the property of quadratics root.

\begin{cor}{Solution of second order homogeneous linear differential equations with constant coefficients}{}
  Let \((aD^2 + bD + c\iden)y = 0\) and let \(r_1, r_2\) be the roots to the polynomial
  \(ax^2 + bx + c\).
  \begin{itemize}
    \item If \(b^2 - 4ac = 0\), then \(r_1 = r_2 = r\) and the solution set is spanned
          by \(x \e^{rx}\) and \(\e^{rx}\).
    \item If \(b^2 - 4ac > 0\), then \(r_1 \neq r_2\) but both of the roots are real;
          hence, the solution set is spanned by \(\e^{r_1x}\) and \(\e^{r_2x}\).
    \item If \(b^2 - 4ac < 0\), then both \(r_1\) and \(r_2\) are complex number of
          the form \(\alpha \pm \im \beta\) where \(\alpha = \frac{b}{2a}\) and
          \(\beta = \sqrt{4ac - b^2}\). The real solution set is then spanned by
          \(\e^{\alpha x} \cos \beta x\) and \(\e^{\alpha x} \sin \beta x\).
  \end{itemize}
\end{cor}

\begin{proof}
  The first two cases are direct consequences of the theorem above.

  For the third one, we can use the quadratic formula to get
  \begin{equation}
    aD^2 + bD + c\iden = (D - \alpha + \im \beta)(D - \alpha - \im \beta)
  \end{equation}
  Hence, the solution set of the differential equation is then spanned by
  \begin{equation}
    \e^{(\alpha+\im \beta)x} \mathand \e^{(\alpha-\im \beta)x}.
  \end{equation}
  Any solution is then a sum of these two functions:
  \begin{align}
    & A \e^{(\alpha+\im \beta)x} + B \e^{(\alpha-\im \beta)x} \nonumber                                                 \\
    & = \e^{\alpha x} (A \e^{\im \beta x} + B \e^{-\im \beta x})                                                        \\
    & = \e^{\alpha x} \ab\big(A \cos \beta x + B \cos (- \beta x) + \im (A \sin \beta x + B \sin (-\beta x))) \nonumber \\
    & = \e^{\alpha x} \ab\big((A + B)\cos \beta x + \im (A - B) \sin \beta x)
  \end{align}
  and we may change the coordinate by letting \(C_1 = A + B\) and \(C_2 = \im (A - B)\).
  Then, the general solution is then
  \begin{equation}
    C_1\e^{\alpha x} \cos \beta x + C_2\e^{\alpha x} \sin \beta x
  \end{equation}
  If \(C_1, C_2 \in \mathbb{R}\), the solution is real and thus
  \begin{equation}
    \e^{\alpha x} \cos \beta x \mathand \e^{\alpha x} \sin\beta x
  \end{equation}
  spans the real solution set.
\end{proof}

Now what if \(P(D)\) has variable coefficients? What can we say about the solution set. Of
course, the kernel of \(P(D)\) still has dimension \(\deg P\). But now, we lost the commutativity.
